\subsection{零 $y^2$ 項の全正常一次境界}
\label{subsec:cubic-third-root-unit-boundary-w0}
\begin{theorem}[零 $y^2$ 項の正常一次chartも空]
\label{thm:cubic-third-root-unit-boundary-w0-empty}
六係数の正常一次chartの $w=0$ 部分も、全係数でH3を与えない。
従って両方向曲率が非零である比例結合の六係数chart全体がH3から除外される。
一方の結合だけが零である層の被覆を、この定理に含めない。
\end{theorem}
\begin{proof}
境界の三関係は
\[
 Dyb+Eby+c^2=0,\quad Eyc+Dcy+b^2=0,\quad Dbc+Ecb=0.
\]
$D,E$ は同時には零でない。
$D\ne0$ なら既知のscalar変更で $D=1$ とし、$E=e$ と書く。
係数体 $\mathbb Q(\omega)(e)$ 上で全次数3,4関係を挿入すると、
商次元は10,15、全関係pivotの非定数因子は $e(e^6-1)$ に限る。
中山固定三次4個の正常行列は $14\times4$。
行ごとの分母を払った多項式行列のSmith因子は
\[
 1,1,1,(e^3-1)(e^3+\omega^2).
\]
左右変換の行列式は非零定数で、全変換を元行列へ掛け戻す。
分母と関係pivotの因子支持も、$e(e^6-1)$ との反復gcdで検証する。
従って $e\ne0,e^6\ne1$ では、正常核が非零となり得るのは
$e^3=-\omega^2$ だけである。

開集合から除かれた値は新しく元関係から検算する。
$e=\pm\omega^j$、$j=0,1,2$ の六値は全て正常核零であり、
$D=0,E=1$ と $D=1,E=0$ も正常核零である。
後二者では次数4商次元は17となるが、実際の商上で正常核を計算している。
各点の非零 $4\times4$ minorを証明書に保持する。
従って全値で、残り得るのは
\[
 D=1,\qquad E^3=-\omega^2.
\tag{W0}
\]

(W0)を $K=\mathbb Q[t]/(t^6-t^3+1)$、$E=t,\omega=t^3-1$ で扱う。
固定した $\omega$ を保つ三埋込みは、(W0)の三点を全て覆う。
$1+\omega E\ne0$ なので既知のshearにより $Z=0$ としてよい。
全 $u,v$ を残し、CY untwist $C$ の元関係から次数3,4を構成する。
正常行列 $27\times8$ の7個のパラメータ定数単元pivotで全残余が零となり、
正常三次は全特殊化で一本の $G$ に限る。
その自由座標を1として、全語係数を公開証明書に記載する。
四正常式への再挿入も多項式として零である。

中心三次は8列の四交換子行列を6個の定数単元pivotで整理して得る。
核は自由階数2であり、$x^3$ から境界で独立な生成元を $\Psi$ とする。
$A=C/(G)$ の重み零中心九次行列は46列、41個の定数単元pivotで全残余零、
自由階数5である。
そのうち $x=0$ で $\Psi^3$ から独立な元を $z_9$ として選ぶ。
$G,\Psi,z_9$ の $x=0$ 像は、全て $u,v$ に依存しない。
全中心核の元四交換子と、全元関係の左右作用を再挿入する。

この選択で、$A/(x,\Psi,z_9)$ のHilbert列は再び
\[
 1,3,6,8,9,9,9,9,9,8,6,3,1,0
\]
となる。H3なら $H_A=(1-t^3)/(1-t)^4$ なので、前節と同じ
有限生成・代数的独立性・Hilbert級数の比較から、
$A$ は $k[x,\Psi,z_9]$ 上自由階数81となる。
これはH3の必要条件からの議論であり、有限次数だけから全次数AS性を主張しない。

全パラメータ共通の境界fiber $A/(x,\Psi-1,z_9-2)$ を使う。
$p=109,t=4,\omega=63$ で補助中央変数 $h$ を使い斉次化すると、
次数12,13は81次元で、$h$ の作用は可逆、
$h=1$ の正則trace行列式は $1\pmod {109}$、中心交換子rankは80である。
元関係の81次行列への再挿入も零である。
正のminorと行列式は標数0に持ち上がるので、このfiberはCSA次数9である。
従って全 $u,v$ に対し、$A$ の中心一般点もCSA次数9、
$Z(A)=k[x,\Psi,z_9]$、$A$ は素である。

$U=(A[x^{-1}])_0$ の全中心は多項式環 $k[\psi,z]$、PI次数9であり、
正常商の双対公式から $\operatorname{Hom}_{k[\psi,z]}(U,k[\psi,z])=U_\delta$、
$\delta^3=1$ となる。
前節で既に証明した局所order・曲面周期の矛盾は、この三つの条件だけを使う。
従ってH3に必要な整域性と左右有限大域次元は両立せず、(W0)も排除される。

公開再現器は \texttt{verify\_cubic\_third\_root\_unit\_boundary\_w0.py}、
同じ接頭辞の \texttt{w0\_center.py} と \texttt{w0\_CSA.py}。
Smith変換、例外値のfresh minor、正常三次・中心元の全係数、
Artinian基底と正のCSA witnessを再生成する。
非零 $w$ の前節と合わせて六係数chart全体の排除を得る。
\end{proof}
